
	This is an interactive problem.
    Your submission will be run against an \emph{interactor},
    which reads from the standard output of your submission
    and writes to the standard input of your submission.
    This interaction needs to follow a specific protocol:

	The interactor first sends one line with an integer $t$ ($1 \leq t \leq 10^4$), the number of test cases.

    For each test case, you can ask \emph{at most two} queries.

    For every query, output ``\texttt{?~$x$~$y$}'' ($0 \leq \left| x \right|$,~$\left| y \right| \leq 50, (x, y) \neq (0, 0)$), where $x$ and $y$ are integers, and $(x, y)$ is the direction you send a ship to.

	The interactor will respond to each query with one line containing a real value $d$ ($0 < d \leq 50$), the distance the ship reported.
    This distance $d$ is given with exactly $10$ decimal places, and has an absolute \emph{and} relative error of at most $10^{-7}$.

    After up to two queries, you should output ``\texttt{!~A}'', where $A$ is an integer, the area of the square.
	It is guaranteed that the square always has a positive integer area.

	The interactor is not adaptive: the borders of the world are fixed before the first query.

    Make sure you \textbf{flush} the standard output after every output.
    For example, you can use \texttt{fflush(stdout)} in C++, \texttt{System.out.flush()} in Java, \texttt{sys.stdout.flush()} in Python, \texttt{std::io::stdout().flush()} in Rust, and \texttt{hFlush stdout} in Haskell.

    A \emph{testing tool} is provided to help you develop your solution.
